\section{Every primitive of the two native presentations}
\label{sec:rules}

The logical rules are shared. The numerical computation schemas differ,
and that difference is kept explicit.

\subsection{Numerical syntax}
The new arithmetic basis has finite terms
\[
 t::=x\mid0\mid St\mid d(t,t)\mid B(n;u,v.h).
\]
Only $u,v$ in $h$ are bound. Its remaining variables are passive parameters,
fixed during a recursion. Put $1=S0$. The two recursion schemas are
\begin{align}
 B(0;u,v.h)&=S h(0,0),\tag{B0}\label{basis:b0}\\
 B(Sn;u,v.h)&=h(n,B(n;u,v.h)).\tag{BS}\label{basis:bs}
\end{align}
The monus schemas are
\[
 d(x,0)=x,\qquad d(x,Sy)=P(d(x,y)),
 \tag{M0, MS}
\]
where
\[
 C(n,a)=B(n;u,v.a),\quad W(n)=B(n;u,v.u),\quad
 P(n)=W(C(n,n)).
\]
All substitutions are simultaneous where appropriate and avoid capture.

The independently measured checker keeps the inherited language
\[
 t::=x\mid0\mid St\mid I(n,a;v.h)
\]
and its equations
\[
 I(0,a;v.h)=a,\qquad I(Sn,a;v.h)=h(I(n,a;v.h)).
 \tag{I0, IS}
\]
The update receives its previous state, not the recursion counter. Its
direct computation packets have the certified-zero syntactic restrictions
specified in Section~\ref{sec:checker}; equality transport restores all
unrestricted instances. These $I$ equations are not axioms of the new
$B+d$ basis.

In both presentations, the binder notation abbreviates a staged ordinary
first-order signature. Each fixed finite earlier-stage body introduces
one function symbol whose arguments are its counter and passive
parameters, together with the independent initializer in the $I$ case.
Bodies are never quantified over. Every term and proof uses finitely many
stages. There is no internal uniformly total evaluator of arbitrary
primitive-recursive codes.

\subsection{Formula constructors and readings}
Native formulas are generated by numerical equations and the paired
binder:
\[
 A::=(s=t)\mid D_{x,y}A.
\]
The auxiliary closure $Q_yA$ may be used in certificates and contexts,
or admitted as a definable formula constructor. Its representative is
$D_{a,y}A$ with unused $a$. Their ordinary readings are
\[
 Q_yA(y):=A(0)\to\forall y\,A(y),\qquad
 D_{x,y}A(x,y):=\exists x\bigl(A(x,0)\to\forall y\,A(x,y)\bigr).
\]
These readings explain the interpretation; ordinary logical rules are
not silently included among the native rules. The designated native
contradiction is $0=1$.

Contexts $\Gamma$ are finite reusable lists. All logical rules are:
\begin{align*}
 &\frac{}{\Gamma\vdash A}\ (A\in\Gamma) &&\text{(Hyp)},\\[4pt]
 &\frac{\Gamma\vdash s=t\quad\Gamma\vdash F(s)}
        {\Gamma\vdash F(t)} &&\text{(Tr)},\\[4pt]
 &\frac{\Gamma\vdash Q_yA(y)\quad\Gamma\vdash A(0)}
        {\Gamma\vdash A(t)} &&\text{(QE)},\\[4pt]
 &\frac{\Gamma\vdash D_{x,y}A(x,y)\quad
         \Gamma,Q_yA(b,y)\vdash F}
        {\Gamma\vdash F} &&\text{(DE)},\\[4pt]
 &\frac{\Gamma,A(t,0)\vdash A(t,Sz)}
        {\Gamma\vdash D_{x,y}A(x,y)} &&\text{(DI-case)},\\[4pt]
 &\frac{\Gamma,E(t,z)\vdash E(t,Sz)}
        {\Gamma\vdash D_{x,y}E(x,y)} &&\text{(DI-progress)}.
\end{align*}
In Tr, $F$ is any fixed native formula template. In DE, $b$ is fresh for
the retained context, conclusion, and fixed matrix parameters. In both
introduction rules, $z$ is fresh for the context, witness, and passive
matrix parameters; the displayed assumption is discharged.
\textbf{DI-progress is restricted to a literal numerical equation.}
An equation beneath a binder does not satisfy that restriction.

The numerical computation schemas are axiom leaves. Equivalently, the
implementation uses
\[
 \frac{\Gamma,H\vdash A}{\Gamma\vdash A}
 \qquad(H\text{ a permitted computation equation}).\tag{Supply}
\]
There are no additional native logical axioms, successor axioms, equality
axioms, or induction schemas. In particular, reflexivity and explosion
are derived below, not omitted costs hidden in a conventional logic.

Weakening, exchange, contraction, renaming, numerical substitution, and
cut of an ordinary assumption are admissible. Their proofs are induction
on finite derivations; eigenparameters are renamed before substitution
or cut insertion. A numerical substitution leaves a literal equation
literal. For syntactically restricted $I$ computations, the certified-zero
repair is applied whenever substitution creates a forbidden instance.

\subsection{The ordinary comparison theory, with its entire logic}
The next section first proves the arithmetic result in a conventional
presentation $T_{Bd}$, to make the invariant easy to inspect. Its
formula grammar is
\[
 A::=(s=t)\mid\bot\mid A\land A\mid A\lor A\mid A\to A
          \mid\forall xA\mid\exists xA.
\]
Its logical rules are exactly intuitionistic natural deduction:

\begin{center}\small
\begin{tabularx}{\linewidth}{@{}lX@{}}
\toprule
Form & Rules\\\midrule
Assumptions & Use any member of the reusable context; weakening, exchange,
and contraction.\\[3pt]
$\bot$ & From $\bot$ infer any formula; no introduction rule.\\[3pt]
$A\land B$ & Introduce from both conjuncts; eliminate to either conjunct.\\[3pt]
$A\lor B$ & Introduce from either disjunct; eliminate by two derivations of
the same conclusion, under the respective disjuncts, discharging them.\\[3pt]
$A\to B$ & Introduce by discharging assumption $A$ from a derivation of $B$;
eliminate with a proof of $A$.\\[3pt]
$\forall xA$ & Introduce from $A(z)$ with fresh $z$ absent from the retained
context and fixed parameters; eliminate by capture-free term substitution.\\[3pt]
$\exists xA$ & Introduce from $A(t)$; eliminate under $A(z)$ with $z$ fresh
for the context, conclusion, and fixed parameters, discharging $A(z)$.\\[3pt]
$=$ & Reflexivity $t=t$; from $s=t$ and $F(s)$ infer $F(t)$ for every
fixed formula template, with capture avoidance.\\
\bottomrule
\end{tabularx}
\end{center}

Its additional arithmetic principles are exactly the four displayed
computation schemas, coverage, separation, and literal-equation induction
listed next. The later native bootstrap derives the conventional logical
interface rather than presupposing it.

Successor is eliminable from the native numerical signature:
$S(a)=C(0,a)$. One can use this abbreviation everywhere, including the
four computation schemas. Keeping $S$ visible makes the proofs easier
to read. This elimination is a definitional equivalence, not an additional
measured checker reduction.
